% Discussions chapter body. Included by thesis.tex under \chapter{DISCUSSIONS}.
% This is an \input fragment, not a standalone document.

\placeholder{Forthcoming. Interpretation of the results, reading the framework comparison across the economic, process, functional, and perception dimensions against the structure and design of the two frameworks. The subsections below are themed to the deeper questions each results section raises.}

\section{Two Realisations of Spec-Driven Development}
\label{sec:disc-sdd}

Spec-kit's planning phase cost more than APM's in ten of the twelve pairs, the one part of the task where APM was the cheaper of the two, as the planning row of Table~\ref{tab:cost-phase} records. The difference follows from how each framework carries out the planning of spec-driven development. Spec-kit plans more autonomously, having the agent explore the codebase and, on ambiguities, decide on best practices, drawing on the model's capabilities and training data. It fills a fixed set of template documents from what it finds, putting few or no questions to the user. APM explores the codebase as well, but it leans more on the user's input and builds review checkpoints into its planning. Required rounds of questions keep the user engaged and abreast of the planning. From what they supply, APM writes three shorter documents that generalise across project types. The two approaches are set out where each framework is introduced, Spec-kit in Section~\ref{sec:frameworks-under-comparison} and APM in Chapter~\ref{chap:framework-design}.

The cost gap has two drivers visible in the planning phase, both pointing the same way, set out in Table~\ref{tab:planning-behaviour}. Spec-kit issued more requests to the model in every one of the twelve pairs, around thirty-nine more at the median and near fifty-seven by the Hodges-Lehmann estimate, with a few pairs where the gap ran much wider. These are the extra steps of an agent that explores and drafts without waiting on the user. Spec-kit also produced more output in nine of the twelve pairs, about thirty-four thousand tokens more at the median and near thirty-six thousand by the Hodges-Lehmann estimate, the text that fills its larger set of documents. APM ran lighter on both.

\begin{table}[H]
  \centering
  \caption{Paired planning-phase request and output differences, APM minus Spec-kit, at 95.8\% coverage.}
  \label{tab:planning-behaviour}
  \begin{tabular}{lrrrrrr}
    \toprule
    Measure & $n$ & Sign ($+$/$-$) & Median & HL & Mean & CI \\
    \midrule
    Requests & 12 & 0 / 12 & $-38.5$ & $-56.5$ & $-60.7$ & $[-95.0, -29.0]$ \\
    Output (thousand tokens) & 12 & 3 / 9 & $-34.5$ & $-35.7$ & $-39.7$ & $[-70.5, -9.9]$ \\
    \bottomrule
  \end{tabular}
\end{table}

The difference is larger under Opus and smaller under Sonnet, yet points the same way in both model groups, as the per-model split in Table~\ref{tab:planning-behaviour-model} of Appendix~\ref{app:sdd-breakdowns} sets out. Spec-kit's heavier request count holds in every pair under both, and its larger output is present in both as well. The typical sizes widen under Opus. With six participants in each group the per-model intervals are wide, so the pooled comparison carries the statistical weight.

In implementation the two frameworks converge. Both build against the planning documents as the source of truth and revise them as findings require, differing in the detail but not the essence. They part on cost. It is there that APM costs more, and that excess is the larger part of the overall difference, coming from how APM splits the work across several agents rather than from how it plans. Section~\ref{sec:disc-cost-mechanics} takes that up.

\placeholder{Forthcoming. The process reading of this realisation: how Spec-kit's autonomous planning shows in the planning phase's turn count and turn time, the more self-driven turns taken over a longer span, set against the lighter user-involved planning of APM.}

\placeholder{Forthcoming. The functional reading of this realisation: the failures connected to how each framework planned, read by their size and type.}

\placeholder{Forthcoming. The perception reading of this realisation, how the two planning approaches relate to what participants reported.}

\section{Cost Mechanics of APM and Spec-kit}
\label{sec:disc-cost-mechanics}

\placeholder{Forthcoming. The implementation-phase cost of each framework and what drives it, including how APM distributes the work across several agents and why the two models diverge.}

\section{Learning Effects}
\label{sec:disc-learning-effects}

The confined number of sessions in the case study allows each participant's pair to be examined individually for learning carried from the first session into the second. The workspace and the agent are reset between sessions, so the participant is the only carrier of that memory, and the evidence is sought in their own inputs: a second-session instruction that pre-empts an error the first session made, an explicit reference to an earlier decision, or a detail handled correctly from recall rather than rediscovered. A stronger second-session result alone does not count, since the agent may have produced it unaided. These signs are read case by case, as indications of learning rather than a conclusive measure of it.

\placeholder{Forthcoming. The per-participant verdicts on carried-over learning and the cohort read drawn from them.}

\section{Threats to Validity}
\label{sec:threats}

This section enumerates the internal, external, and construct validity threats and the mitigations applied. Each threat points back to the specific section that motivates it.

\subsection{Internal Validity}

\textbf{Learning effects from the within-subjects design.} The within-subjects design declared in the Case-Study Design (Section~\ref{sec:case-study-design}) has each participant complete the same task twice, so task familiarity from the first session can influence performance in the second. Counterbalancing the framework order across participants, as described in Sessions and Data (Section~\ref{sec:sessions-data}), mitigates the framework-effect confound but does not eliminate the underlying learning curve. The learning-effects examination described in Section~\ref{sec:disc-learning-effects} reads each participant's two sessions for evidence of carried-over learning and is the primary defence against this threat.

\textbf{Self-reported versus observed measures.} The perception measures described in Section~\ref{sec:measures} pair the instrumented, transcript-derived measures with self-reported workload, confidence, and perceived framework effectiveness; observed divergences between self-report and instrument are reported as findings rather than absorbed as an unstated limitation. Residual risk remains: self-reports are subject to recall bias and social-desirability pressures that the instrument-side measurements cannot independently correct.

\subsection{External Validity}

\textbf{Generalisation scope.} The case study is run on a single AI coding CLI, against a complex software engineering task requiring two-language implementation and legacy-codebase research, described in the Task Specification (Section~\ref{sec:task-spec}), in a setting positioned between laboratory and field conditions. Findings generalise within this scope and do not generalise to multi-platform settings, dynamic cross-model routing, simpler or less knowledge-intensive tasks, or pure laboratory or natural-setting conditions.

\textbf{Temporal decay of findings.} The AI coding assistant field is evolving rapidly. Effect sizes and descriptive claims throughout this chapter (rate-card prices, command sets, framework releases, and model availability) are tied to the May 2026 snapshot stated in Section~\ref{sec:methodology}. The dollar-equivalent cost values reported under Chapter~\ref{chap:results} rest on the rate-card snapshot; results that rest on those specific rates are defensible against the snapshot but may require re-estimation under a materially different rate card.

\subsection{Construct Validity}

\textbf{Behavioural input buckets as process-efficiency proxy.} The short, protocol, and standard input buckets defined under the Process measures in Section~\ref{sec:measures} are extracted from the transcript content for this study. They are not validated against an external coordination-burden instrument. Their validity as proxies for human coordination effort rests on face validity alone: the relative ordering of frameworks along these metrics is robust, but the absolute magnitude of differences is interpretive rather than calibrated.

\textbf{Instrument coverage limits.} The 100-point rubric captures correctness and robustness but not code quality, maintainability, security posture, or long-term technical debt. Metric extraction is further limited to the Claude Code JSONL transcripts and their sibling subagent files that accompany each participant submission, as declared in Sessions and Data (Section~\ref{sec:sessions-data}): any interaction with the AI that took place outside the CLI's transcript surface is invisible to the metric set.

\textbf{Replica versus live Open eClass.} The shared Docker-based replica of the Open eClass platform described in the Task Environment (Section~\ref{sec:task-env}) pins the platform to a fixed upstream baseline and seeds deterministic fixture data. The shared replica makes the cohort comparable and the grading reproducible, and removes any parity gap between what participants build against and what the evaluation runs against. The trade-off is that the task is performed against a replica rather than against any live university instance of Open eClass: any property of a live instance that the replica does not reproduce, such as timezone handling of locale-formatted dates, session-cookie behaviour under load, cross-origin interactions with the mock CAS service, or upstream behaviours that have shifted since the baseline commit, is not exercised by the case study.
